% refs.pdf: smart references, outline, page labels, links and search normalization.
% Facts the tests rely on are listed in ../README.md; keep them in sync.
\pdfinfoomitdate=1
\pdftrailerid{}
\pdfsuppressptexinfo=-1
\documentclass[11pt]{article}
\pdfpagewidth=\paperwidth \pdfpageheight=\paperheight % US letter whatever the local TeX default
\usepackage[T1]{fontenc}
\usepackage{amsmath,amsthm}
\usepackage{cite}
\renewcommand{\citepunct}{,\ }
\usepackage[pdftitle={Folio Test Paper},pdfauthor={Folio Tests},pdfcreator={},pdfproducer={}]{hyperref}
\numberwithin{equation}{section}
\newtheorem{theorem}{Theorem}[section]
\newtheorem{lemma}{Lemma}[section]
\begin{document}

\pagenumbering{gobble}
\begin{center}
  {\LARGE Folio Test Paper}\\[1em]
  Title page without a page number
\end{center}
\clearpage

\pagenumbering{roman}
\section*{Preface}
This front matter is numbered with roman numerals.
\clearpage
\noindent The second page of front matter.
\clearpage

\pagenumbering{arabic}
\section{Introduction}
\label{sec:intro}
Early work on the subject is reviewed in the survey \cite{alpha}, and the classic
results are due to two groups \cite{beta,gamma}. Several later papers extend this \cite{delta,epsilon,zeta}.
The data are collected in Table~\ref{tab:data}, and the setup is drawn in Figure~\ref{fig:setup}.
The main result is Theorem~\ref{thm:main}, and the key estimate is Lemma~\ref{lem:key}.

\subsection{Motivation}
Every field of physics has its ``quoted'' folklore -- and its exceptions---like this one.
The word renormaliza-\newline tion is broken across two lines on purpose.
Details are in Section~\ref{sec:results}. More is available at
\url{https://example.org/folio} and in \href{file:notes.txt}{the local notes}.

\begin{theorem}
\label{thm:first}
Every bounded sequence has a convergent subsequence.
\end{theorem}

\begin{theorem}
\label{thm:main}
The main estimate holds for all inputs.
\end{theorem}

\begin{figure}[h]
  \centering
  \fbox{\rule{0pt}{2cm}\hspace{4cm}}
  \caption{The experimental setup.}
  \label{fig:setup}
\end{figure}
\clearpage

\section{Results}
\label{sec:results}
The basic relation is
\begin{equation}
  \label{eq:basic}
  E = m c^2
\end{equation}
and a second relation follows:
\begin{equation}
  \label{eq:second}
  F = m a .
\end{equation}
Combining Eq.~\eqref{eq:basic} with the second one gives the claim, and
the relation (2.1) is used again below.

\begin{table}[h]
  \centering
  \caption{A first table.}
  \begin{tabular}{cc} a & b \\ c & d \end{tabular}
\end{table}

\begin{table}[h]
  \centering
  \caption{Measured data.}
  \label{tab:data}
  \begin{tabular}{cc} 1 & 2 \\ 3 & 4 \end{tabular}
\end{table}
\clearpage

\section{Proofs}
\begin{lemma}
\label{lem:key}
The key estimate is sharp.
\end{lemma}
\begin{proof}
Use Theorem~\ref{thm:first} and the relation in Eq.~\eqref{eq:second}.
\end{proof}
\clearpage

\begin{thebibliography}{9}
\bibitem{alpha} A. Alpha, A survey of the field, J. Tests 1 (2020).
\bibitem{beta} B. Beta, Classic results, J. Tests 2 (2021).
\bibitem{gamma} C. Gamma, More classic results, J. Tests 3 (2021).
\bibitem{delta} D. Delta, An extension, J. Tests 4 (2022).
\bibitem{epsilon} E. Epsilon, Another extension, J. Tests 5 (2023).
\bibitem{zeta} Z. Zeta, The last extension, J. Tests 6 (2024).
\end{thebibliography}
\end{document}
